% Chapter 1

\chapter{Conventions} % Main chapter title

\label{Chapter2} % For referencing the chapter elsewhere, use \ref{Chapter1}

%----------------------------------------------------------------------------------------

% Define some commands to keep the formatting separated from the content
\renewcommand{\keyword}[1]{\textbf{#1}}
\renewcommand{\tabhead}[1]{\textbf{#1}}
\renewcommand{\code}[1]{\texttt{#1}}
\renewcommand{\file}[1]{\texttt{#1}}
\renewcommand{\option}[1]{\texttt{\itshape#1}}

%----------------------------------------------------------------------------------------



An important property of a long document such as a thesis is consistency.%
\todo{This is a todo note. Use them to ensure consistency. When you print the final thesis, hide all todo notes by adding the parameter \emph{final} to the documentclass declaration.}
It is, however, quite difficult to keep track of terminology, style, and structure over an extended period of time.

Following certain conventions can help you to achieve a consistent result. This chapter summarizes conventions and recommendations for effective thesis writing.\marginnote{Having a signpost paragraph before the first section in a chapter is considered good style.}

This chapter contains only selected recommendations. For a more comprehensive list that also includes links to useful tools, consider \url{https://github.com/TheHairyJ/StudentResources/blob/master/Writing.md}.

\section{Links}

Special care is necessary when you include links in your thesis. URLs are generally embedded using the \verb|\url{}| command. This will allow \LaTeX{} to break them at the end of a line. Moreover, they will become clickable.

It is an open question whether \emph{all} links should be embedded as proper references or not. We recommend to follow these principles:

\begin{itemize}
  \item \textbf{Use a link} in the text (or as a footnote when it is a longer link) when the main purpose of providing the URL is to \emph{provide a reference to that particular location}. Example: links to the website of a service or product or to the homepage of a weblog, when the purpose is to make readers aware of its existence as a whole.
  \item \textbf{Use a reference} when the purpose is to \emph{refer to or discuss the content} on a website. References should generally mention the author and the publication date (at least the year). If the author cannot be identified, you can resort to using the name of the respective organization instead (use an additional pair of curly braces in BibTeX -- otherwise the name will be mis-interpreted as first name and last name). Example: linking to a particular blog post.
\end{itemize}

\paragraph{Link Rot}

\marginnote{See also: \href{https://en.wikipedia.org/wiki/Link_rot}{Link Rot in Wikipedia}}As the web is quite dynamic, many links do not work forever. Content may move to different URLs, it may be substantially altered, and sites may go out of service altogether. This phenomenon is called Link Rot, and it has become a significant issue in academia.

We generally recommend to \emph{archive all websites} whose URLs you provide using \textsc{The Wayback Machine} of the Internet Archive, which is available at \url{archive.org}.
This preserves the content of a website at the given point in time and makes it permanently available in the Internet Archive.

\marginnote{All links in this document have been archived at the time of its publication.}You should inform the reader about the availability of an archived version. We recommend to provide the \emph{original link} and add ``(archived at [date])'' after the URL. This applies to URLs in the main text, in footnotes, and in references.

Archiving is not necessary for academic papers that have a Digital Object Identifier (DOI), have an ISBN/ISSN, or are available on \url{arxiv.org}.


\section{References}

Scientific references should come \emph{before} the punctuation mark (such as a comma or period) if there is one, as shown in this sentence \cite{Hintz02}.

When you cite multiple references in one spot, you should include them in one \verb|\cite| command:
\begin{quote}
  The FOO approach appears in a wide variety of applications \cite{murdoch_steven_j._chip_2010,anderson_ross_emv:_2014,kou_weidong_secure_2003}.
\end{quote}

Note, however, that this style is not very useful for the reader.\footnote{This recommendation was originally given here: \url{https://nhigham.com/2014/12/22/more-tips-on-book-and-thesis-writing/}} Much more information can be given like this:
\begin{quote}
The FOO approach appears in a wide variety of applications, such as \emph{qux systems} \cite{murdoch_steven_j._chip_2010}, \emph{baz systems} \cite{anderson_ross_emv:_2014}, and -- much earlier -- in \emph{bar readers} \cite{kou_weidong_secure_2003}.
\end{quote}

A remark on style: References should not be treated as nouns. Avoid writing ``\cite{murdoch_steven_j._chip_2010} shows that many EMV systems are broken \ldots'' and replace it with ``Murdoch et al. show that EMV systems are broken  \cite{murdoch_steven_j._chip_2010}.''

\paragraph{BibTeX Snippets}

For many publications you will find BibTeX snippets online, for instance in the ACM Digital Library or on \url{https://dblp.org}. Using these snippets is generally recommended. Unfortunately, the quality of publicly available BibTeX snippets varies a lot. This means that you will have to edit them to create a consistent bibliography. For instance, make sure that the name of a particular conference is always spelled consistently in all entries. Moreover, you should remove superfluous information such as precise dates, multiple occurrences of the year, and the location of the conference.

A particular issue, is capitalization in titles. Whether you use sentence case or title case is up to you -- as long as you stick to the choice. You may have to enclose the whole title of a bibitem in double curly braces to preserve capitalized letters. Alternatively, you can only enclose capitalized words or abbreviations in curly braces. Consider the following example in which we ensure that TCP is printed in uppercase:
\begin{verbatim}
   @inproceedings{... title = "Improving {TCP}-based Communication Systems" ...}
\end{verbatim}

\section{Margin Notes and Footnotes}

The template offers you the option to use margin notes and footnotes. You can use Margin notes (\verb|\marginnote|) as a replacement for footnotes. You can also use them with keywords or short phrases to attract the reader's attention.

Footnotes are either inserted directly after the particular word they refer to or after the punctuation if they relate to the whole clause or sentence.\footnote{Such as this footnote, here down at the bottom of the page: Lorem ipsum dolor sit amet, consectetur adipisicing elit, sed do eiusmod tempor incididunt ut labore et dolore magna aliqua.} Footnotes themselves should be full, descriptive sentences (beginning with a capital letter and ending with a full stop). The APA6 states: \enquote{Footnote numbers should be superscripted, [...], following any punctuation mark except a dash.} The Chicago manual of style states: \enquote{A note number should be placed at the end of a sentence or clause. The number follows any punctuation mark except the dash, which it precedes. It follows a closing parenthesis.}

Using both margin notes and footnotes for similar purposes may create an inconsistent appearance and is therefore discouraged.

\section{Tables}

Tables are an important way of displaying your results. Below is an example table which was generated with this code:

{\small
\begin{verbatim}
\begin{table}
\caption{The effects of treatments X and Y on the four groups studied.}
\label{tab:treatments}
\centering
\begin{tabular}{r r r}
\toprule
& \multicolumn{2}{c}{\tabhead{Results}} \\ \cmidrule(lr){2-3}
\tabhead{Group} & \tabhead{Treatment X} & \tabhead{Treatment Y} \\
\midrule
1 & 0.2 & 0.8\\
2 & 0.17 & 0.7\\
3 & 0.24 & 0.75\\ \addlinespace
4 & 0.68 & 0.3\\
5 & 0.61 & 0.9\\
6 & 0.18 & 0.1\\
\bottomrule\\
\end{tabular}
\end{table}
\end{verbatim}
}

\begin{table}
\caption{The effects of treatments X and Y on the four groups studied.}
\label{tab:treatments}
\centering
\begin{tabular}{r r r}
\toprule
& \multicolumn{2}{c}{\tabhead{Results}} \\ \cmidrule(lr){2-3}
\tabhead{Group} & \tabhead{Treatment X} & \tabhead{Treatment Y} \\
\midrule
1 & 0.2 & 0.8\\
2 & 0.17 & 0.7\\
3 & 0.24 & 0.75\\ \addlinespace
4 & 0.68 & 0.3\\
5 & 0.61 & 0.9\\
6 & 0.18 & 0.1\\
\bottomrule\\
\end{tabular}
\end{table}

Column heads should use sentence-format capitalization (see \url{http://www.chicagomanualofstyle.org/15/ch13/ch13_sec019.html}).

Tables have captions that always appear \emph{above} the \code{tabular} environment. All tables must be referred to in the main text.
You can reference tables with \verb|\ref{<label>}| where the label is defined within the table environment. See \file{Chapter1.tex} for an example of the label and citation (e.\,g., Table~\ref{tab:treatments}).

In the PSI Thesis template tables use the \emph{booktabs} style, which avoids visual clutter such all vertical lines. If you want to group columns, do it in the headers using \verb|\cmidrule|.

Horizontal lines should be used sparingly as well. You can add additional vertical space (using \verb|\addlinespace|) to group consecutive rows (as was done in Table~\ref{tab:treatments}). You can also use bold print to highlight especially relevant parts.

Note that Table~\ref{tab:treatments} is not laid out ideally: the numbers are given with different precision, which disturbs their alignment. Besides, strictly aligning  all numeric columns to the right is not always appropriate.
The effect of changing to centered alignment can be seen in Table~\ref{tab:treatments2}.

\begin{table}
\caption{The effects of treatments X and Y on the four groups studied.}
\label{tab:treatments2}
\centering
\begin{tabular}{c c c}
\toprule
& \multicolumn{2}{c}{\tabhead{Results}} \\ \cmidrule(lr){2-3}
\tabhead{Group} & \tabhead{Treatment X} & \tabhead{Treatment Y} \\
\midrule
1 & 0.20 & 0.80\\
2 & 0.17 & 0.70\\
3 & 0.24 & 0.75\\ \addlinespace
4 & 0.68 & 0.30\\
5 & 0.61 & 0.90\\
6 & 0.18 & 0.10\\
\bottomrule\\
\end{tabular}
\end{table}


Creating good tables is a challenge in its own. Recommended readings: 
\begin{itemize}
\item \url{https://nhigham.com/2019/11/19/better-latex-tables-with-booktabs/}
\item \url{http://tug.org/pracjourn/2007-1/mori/mori.pdf}
\item \url{https://inf.ethz.ch/personal/markusp/teaching/guides/guide-tables.pdf}
\item \url{https://www.cl.uni-heidelberg.de/courses/ss19/wissschreib/material/tableTricks.pdf}
\end{itemize}

If you are in a hurry, you may find \url{https://www.tablesgenerator.com} and \url{https://www.latex-tables.com} useful. You can also use \texttt{LyX} to build tables and export their source code.


\section{Figures}

There will hopefully be many figures in your thesis (that should be placed in the \emph{Figures} folder). The way to insert figures into your thesis is to use a code template like this:
\begin{verbatim}
\begin{figure}[t]
\centering
\includegraphics[width=0.75\textwidth]{Figures/Electron}
\decoRule
\caption[An Electron]{An electron (artist's impression).}
\label{fig:Electron}
\end{figure}
\end{verbatim}
Also look in the source file. Putting this code into the source file produces the picture of the electron that you can see in the figure below.

\begin{figure}[t]
\centering
\includegraphics[width=0.75\textwidth]{Figures/Electron}
\decoRule
\caption[An Electron]{An electron (artist's impression).}
\label{fig:Electron}
\end{figure}

The recommended figure placement is the top of the page (denoted by \code{[t]}). Don't worry about figures not appearing exactly where you write them in the source. The placement depends on how much space there is on the page for the figure. Sometimes there is not enough room to fit a figure directly where it should go (in relation to the text) and so \LaTeX{} puts it at the top of the next page. Positioning figures is the job of \LaTeX{} and so you should only worry about making them look good. Generally, everything is fine as long as a figure either appears on the page where it is referred to for the first time or on one of the subsequent pages.

Every figure needs a descriptive caption and a label. Figure captions must always appear below the included graphics file within the \code{figure} environment.

Every figure must be referred to in the text at least once, either in parentheses (cf. Fig.~\ref{fig:Electron}) or explicitly in the sentence: Figure~\ref{fig:Electron} shows an electron. Refer to figures using the abbreviation \code{Fig.} followed by a protected space (\verb|~|) and \verb|\ref|. Exception: write \code{Figure} if it is the first word of a sentence.

The \verb|\caption| command contains two parts,%
\marginnote{Theses at PSI usually do not have a List of Figures. You can, therefore, omit the part in square brackets.} 
the first part, inside the square brackets is the title that will appear in the \emph{List of Figures}, and so should be short.
 The second part in the curly brackets should contain the longer and more descriptive caption text.

The \verb|\decoRule| command is optional and simply puts a horizontal line below the image. This can be useful with figures that are not symmetrical or whose exterior is not evenly defined at the bottom.

Resize your figures, ideally consistently, to an appropriate size, e.\,g., a fraction of \texttt{\textbackslash textwidth}.

\LaTeX{} is capable of using images in PDF, JPG, and PNG format. Using vectorized figures (PDF) is preferred as they usually provide sharper results. If you have to use pixel graphics, create them with a sufficiently high resolution, i.\,e, at least 300 dpi.

\paragraph{Side-by-Side Figures} Sometimes it is desirable to show multiple figures side-by-side. There are several approaches, for instance, by using the \texttt{subfigure} package. In many cases, a comprehensive subfigure support is not really needed. A lightweight approach as shown in Fig.~\ref{fig:sidebyside} may be sufficient. This result is achieved with the following code:

\begin{verbatim}
\begin{figure}[t]
\centering
\includegraphics[width=0.48\textwidth]{Figures/Electron}%
\hspace{\fill}%
\includegraphics[width=0.48\textwidth]{Figures/Electron}
\decoRule
\caption{A blue electron (left) and another one (right).}
\label{fig:sidebyside}
\end{figure}
\end{verbatim}

\begin{figure}[t]
\centering
\includegraphics[width=0.48\textwidth]{Figures/Electron}%
\hspace{\fill}%
\includegraphics[width=0.48\textwidth]{Figures/Electron}
\decoRule
\caption{A blue electron (left) and another one (right).}
\label{fig:sidebyside}
\end{figure}

\section{Typesetting Mathematics}

If your thesis is going to contain heavy mathematical content, be sure that \LaTeX{} will make it look beautiful, even though it won't be able to solve the equations for you.

The \enquote{Not So Short Introduction to \LaTeX} (available on \href{http://www.ctan.org/tex-archive/info/lshort/english/lshort.pdf}{CTAN}) should tell you everything you need to know for most cases of typesetting mathematics. If you need more information, a much more thorough mathematical guide is available from the AMS called, \enquote{A Short Math Guide to \LaTeX} and can be downloaded from:
\url{ftp://ftp.ams.org/pub/tex/doc/amsmath/short-math-guide.pdf}

There are many different \LaTeX{} symbols to remember, luckily you can find the most common symbols in \href{http://ctan.org/pkg/comprehensive}{The Comprehensive \LaTeX~Symbol List}.

You can write an equation, which is automatically given an equation number by \LaTeX{} like this:
\begin{verbatim}
\begin{equation}
E = mc^{2}
\label{eqn:Einstein}
\end{equation}
\end{verbatim}

This will produce Einstein's famous energy-matter equivalence equation:
\begin{equation}
E = mc^{2}
\label{eqn:Einstein}
\end{equation}

All equations you write are automatically given equation numbers by \LaTeX{}. If you don't want a particular equation numbered, use the unnumbered form:
\begin{verbatim}
\[ a^{2}=4 \]
\end{verbatim}
You can also have equations in the middle of a paragraph: \( E = mc^{2} \). This can be achieved with this code: 
\begin{verbatim}
\( E = mc^{2} \)
\end{verbatim}

%----------------------------------------------------------------------------------------

\section{Structure}

You should break your thesis up into chapters and sections. \LaTeX{} automatically builds a Table of Contents by looking at all the \verb|\chapter{}|, \verb|\section{}|  and \verb|\subsection{}| commands you write in the source. You may even think about using \emph{sub}subsections (\verb|\subsubsection{}|). All of these will be hierarchically numbered.

Avoid dangling elements in the hierarchy. This means that if you have a Section~1.1, you should also have (at least) Section~1.2. This is true for all levels of the hierarchy.

Note, however, that (sub-)subsections should only be used if you use them consistently in all or multiple chapters. Otherwise, you should opt for a \verb|paragraph{}|, which structures pieces of content with unobtrusive non-numbered inline headings.

The Table of Contents should only list \emph{chapters} and \emph{sections}. Listing \emph{(sub-)subsections} is not recommended as this results in a very long Table of Contents, which may become difficult to read. The depth to which the Table of Contents is formatted is set within \file{setup.tex}. If you need this changed, search for \texttt{\textbackslash etocsettocdepth}.

Use \emph{title case} in your \emph{chapters}, \emph{sections}, and \emph{paragraphs}. Consider, for instance, \url{https://capitalizemytitle.com} for title case capitalization rules.


\section{Typography}

Use thin spaces (\verb|\,|) in the appropriate places. For instance, write \verb|i.\,e.,| to obtain i.\,e., which is spoken ``that is''. The same applies to e.\,g. Both abbreviations are followed by a comma in American English.

Dashes can be used instead of colons or pairs of commas to mark off a nested clause. Use an \emph{en dash} for that purpose -- like in this example. If you cannot type an \emph{en dash} on your keyboard, you can write two regular hyphens next to each other. \LaTeX{} will substitute them with an \emph{en dash}. While we have no strict preference, we recommend to use \emph{en dashes} instead of the longer \emph{em dashes}.

Use en dashes also within ranges, e.\,g., when you write something like 5--10\,\% (note the thin space before the \% symbol).

Use a proper \emph{minus symbol} (e.\,g., by using math mode like this \(-1.337\)). Also, use a sensible number of digits after the decimal point.

Use proper typographic quotation marks like ``these'' (which can be obtained with \verb|``these''| -- or by copy-and-pasting the respective Unicode characters).

If you speak German, consider reading \textsc{typokurz} (\url{https://zvisionwelt.wordpress.com/downloads/}), a short introduction into typographic issues. More detailed information can be found on Matthew Butterick's website \url{https://practicaltypography.com}.

\section{Style}

We highly recommend to use the serial comma in all lists to avoid ambiguity. The serial comma is also known as the \emph{Oxford comma}. It is inserted before the word \emph{and} that leads the last item of a list. Consider this example:\footnote{Source: \url{https://nhigham.com/2016/02/16/the-serial-or-oxford-comma/}}
\begin{quote}
  Three important techniques in the design of algorithms are bisection, divide and conquer, and recursion.
\end{quote}

``Lists are common in all forms of writing. The list items can be included within the text or put on separate lines. Separate lines are used in order to draw attention to the items, to ease reading when the items are long or numerous, or to facilitate cross-reference to specific items.''\footnote{Source: \url{https://nhigham.com/2015/12/17/punctuating-lists/}} The environments \code{itemize} and \code{enumerate} produce lists on separate lines. It is considered good style to use punctuation in such a way that the lists remain readable as full sentences if their items were not split onto separate lines. This can be achieved by ensuring that every list item is a full sentence. If list items are fragments, additional punctuation is necessary (example taken from the source mentioned in the footnote):

\begin{quote}
  We used three different algorithms in the experiments. The table reports the performance of
\begin{itemize}
\item Algorithm 3.1 (based on a Taylor series),
\item Algorithm 3.2 (with parameter \(k = 1\)), and
\item Algorithm 3.3 (with tolerance \(10^{-8}\)).
\end{itemize}
\end{quote}

\paragraph{Common Bugs in Writing}

Read the comprehensive list of common bugs in writing by Henning Schulzrinne, which is available at \url{http://www.cs.columbia.edu/~hgs/etc/writing-bugs.html}.